% Chapter 2: First-Order Differential Equations
% Template file - Replace this content with your chapter material

%% Preamble - Common setup for all beamer presentations
%% Include this file in your chapter presentations

\documentclass[aspectratio=169, 12pt]{beamer}

% Packages
\usepackage[utf8]{inputenc}
\usepackage[T1]{fontenc}
\usepackage{amsmath}
\usepackage{amssymb}
\usepackage{amsthm}
\usepackage{mathtools}
\usepackage{graphicx}
\usepackage{xcolor}
\usepackage{tcolorbox}
\usepackage{tikz}
\usetikzlibrary{shadows, arrows, shapes}
\usepackage{hyperref}
\usepackage{listings}
\usepackage{geometry}

% Custom colors (import from custom_colors.tex if available)
\definecolor{defcolor}{RGB}{230, 245, 250}
\definecolor{defborder}{RGB}{0, 102, 204}
\definecolor{thmcolor}{RGB}{245, 230, 250}
\definecolor{thmborder}{RGB}{153, 51, 153}
\definecolor{excolor}{RGB}{240, 255, 240}
\definecolor{exborder}{RGB}{34, 139, 34}
\definecolor{probcolor}{RGB}{255, 250, 240}
\definecolor{probborder}{RGB}{255, 140, 0}

% Beamer Theme
\usetheme{Madrid}
\usecolortheme{default}
\setbeamertemplate{navigation symbols}{}
\setbeamertemplate{footline}[frame number]

% tcolorbox environment definitions

%% Definition Box
\newtcolorbox{definitionbox}[1]{
  colback=defcolor,
  colframe=defborder,
  coltitle=white,
  colbacktitle=defborder,
  title=#1,
  fonttitle=\bfseries\large,
  boxrule=2pt,
  arc=5pt,
  left=8pt,
  right=8pt,
  top=8pt,
  bottom=8pt,
  shadow={2pt}{-2pt}{0pt}{black!20},
  enhanced
}

%% Theorem Box
\newtcolorbox{theorembox}[1]{
  colback=thmcolor,
  colframe=thmborder,
  coltitle=white,
  colbacktitle=thmborder,
  title=#1,
  fonttitle=\bfseries\large,
  boxrule=2pt,
  arc=5pt,
  left=8pt,
  right=8pt,
  top=8pt,
  bottom=8pt,
  shadow={2pt}{-2pt}{0pt}{black!20},
  enhanced
}

%% Example Box
\newtcolorbox{examplebox}[1]{
  colback=excolor,
  colframe=exborder,
  coltitle=white,
  colbacktitle=exborder,
  title=#1,
  fonttitle=\bfseries\large,
  boxrule=2pt,
  arc=5pt,
  left=8pt,
  right=8pt,
  top=8pt,
  bottom=8pt,
  shadow={2pt}{-2pt}{0pt}{black!20},
  enhanced
}

%% Problem Box
\newtcolorbox{problembox}[1]{
  colback=probcolor,
  colframe=probborder,
  coltitle=white,
  colbacktitle=probborder,
  title=#1,
  fonttitle=\bfseries\large,
  boxrule=2pt,
  arc=5pt,
  left=8pt,
  right=8pt,
  top=8pt,
  bottom=8pt,
  shadow={2pt}{-2pt}{0pt}{black!20},
  enhanced
}

% Custom commands
\newcommand{\RR}{\mathbb{R}}
\newcommand{\CC}{\mathbb{C}}
\newcommand{\NN}{\mathbb{N}}
\newcommand{\ZZ}{\mathbb{Z}}
\newcommand{\dd}[1]{\frac{d}{d#1}}
\newcommand{\dydx}{\frac{dy}{dx}}
\newcommand{\dd2ydx2}{\frac{d^2y}{dx^2}}

% Title page formatting
\title[Your Chapter Title]{Your Chapter Title}
\author{Author Name}
\institute{Institution}
\date{\today}

% Remove the footer from title page
\setbeamertemplate{footline}{}


\title[Chapter 2]{Chapter 2: First-Order Differential Equations}
\subtitle{Solution Methods and Applications}
\author{Your Name}
\date{\today}

\begin{document}

\begin{frame}
    \titlepage
\end{frame}

\begin{frame}
    \frametitle{Chapter Overview}
    \tableofcontents
\end{frame}

\section{Separable Equations}

\begin{frame}
    \frametitle{Separable Differential Equations}
    
    \begin{definitionbox}{Separable Equation}
        A first-order differential equation of the form
        \[
        M(x) + N(y)\dydx = 0
        \]
        is called \textbf{separable}.
    \end{definitionbox}
\end{frame}

\section{Linear Equations}

\begin{frame}
    \frametitle{Linear First-Order Equations}
    
    \begin{definitionbox}{Standard Form}
        A linear first-order differential equation has the form:
        \[
        \dydx + P(x)y = Q(x)
        \]
        where $P(x)$ and $Q(x)$ are continuous functions.
    \end{definitionbox}
\end{frame}

\section{Exact Equations}

\begin{frame}
    \frametitle{Exact Differential Equations}
    
    \begin{theorembox}{Exactness Criterion}
        The equation $M(x,y) + N(x,y)\dydx = 0$ is exact if and only if:
        \[
        \frac{\partial M}{\partial y} = \frac{\partial N}{\partial x}
        \]
    \end{theorembox}
\end{frame}

\section{Applications}

\begin{frame}
    \frametitle{Applications of First-Order ODE}
    
    \begin{itemize}
        \item Exponential growth and decay
        \item Newton's law of cooling
        \item Population dynamics
        \item Chemical reactions
    \end{itemize}
\end{frame}

\begin{frame}
    \frametitle{Example: Exponential Decay}
    
    \begin{problembox}{Problem}
        A radioactive substance has a half-life of 1600 years. If there are initially 100 grams, find the amount remaining after 3200 years.
    \end{problembox}
\end{frame}

\begin{frame}
    \frametitle{Example Solution}
    
    \begin{examplebox}{Solution}
        The exponential decay model is $\frac{dA}{dt} = -kA$.
        
        Solution: $A(t) = A_0 e^{-kt}$
        
        Using half-life: $50 = 100e^{-k \cdot 1600}$
        
        Solving: $k = \frac{\ln 2}{1600}$
        
        After 3200 years: $A(3200) = 100 e^{-k \cdot 3200} = 100 e^{-2\ln 2} = 25$ grams
    \end{examplebox}
\end{frame}

\section{Summary}

\begin{frame}
    \frametitle{Summary}
    
    \begin{itemize}
        \item Separable equations can be solved by integration
        \item Linear equations use integrating factor method
        \item Exact equations have special solution structure
        \item Many real-world phenomena follow first-order models
    \end{itemize}
\end{frame}

\end{document}
